\section{Setup}
\label{sec:setup}

\paragraph{Scope.}
A developer evaluates a chat model at full precision and ships a 4-bit copy, and we ask whether the copy refuses the same requests. We use assistant prefill, which forces the first tokens of the reply, to test how easily refusal is bypassed \citep{qi2024shallow}. It is not a realistic attack on local weights, whose holder can remove safety training directly.

\paragraph{Models and formats.}
We use \qwen-Instruct, \llama-Instruct and \qwenthree-Instruct-2507 with their default chat templates, pinning Llama's template date so that runs are comparable. Every quantized build is converted from the bf16 checkpoint with \texttt{mlx-lm} 0.32.2. Our primary format is the \texttt{mlx-lm} default, affine round-to-nearest at 4 bits with groups of 64. Around it we build 8 bits, 4 bits with groups of 32, \texttt{mlx-lm}'s mixed 4/6-bit recipe and MXFP4. We also build \texttt{mlx-lm}'s AWQ and GPTQ \citep{lin2024awq, frantar2023gptq}, which write the same 4-bit, group-64 format as the default, so their bits per weight are the same. For a second runtime we build bf16, Q8\_0 and Q4\_K\_M GGUF files with llama.cpp from the same checkpoints and send it the same prompts as token ids. Table~\ref{tab:formats} gives each format's bits per weight and its capability loss, measured as the change in negative log-likelihood (NLL) that the build assigns to bf16's own replies to 200 harmless prompts.

\begin{table}[t]
\centering\small
\setlength{\tabcolsep}{3.5pt}
\begin{tabular}{lcccc}
\toprule
Format & bits/w & \multicolumn{3}{c}{$\Delta$NLL/token vs bf16} \\
 & & Qwen2.5 & Llama & Qwen3 \\
\midrule
8bit-g64   & \N{cap.qwen.8bit-g64.bpw}  & \N{cap.qwen.8bit-g64.dnll}  & \N{cap.llama.8bit-g64.dnll}  & \N{cap.qwenthree.8bit-g64.dnll} \\
4bit-g32   & \N{cap.qwen.4bit-g32.bpw}  & \N{cap.qwen.4bit-g32.dnll}  & \N{cap.llama.4bit-g32.dnll}  & \N{cap.qwenthree.4bit-g32.dnll} \\
mixed\_4\_6 & \N{cap.qwen.mixed46.bpw} & \N{cap.qwen.mixed46.dnll}   & \N{cap.llama.mixed46.dnll}   & \N{cap.qwenthree.mixed46.dnll} \\
4bit-g64   & \N{cap.qwen.4bit-g64.bpw}  & \N{cap.qwen.4bit-g64.dnll}  & \N{cap.llama.4bit-g64.dnll}  & \N{cap.qwenthree.4bit-g64.dnll} \\
MXFP4      & \N{cap.qwen.mxfp4.bpw}     & \N{cap.qwen.mxfp4.dnll}     & \N{cap.llama.mxfp4.dnll}     & \N{cap.qwenthree.mxfp4.dnll} \\
AWQ        & \N{cap.qwen.awq.bpw}       & \N{cap.qwen.awq.dnll}       & \N{cap.llama.awq.dnll}       & -- \\
GPTQ       & \N{cap.qwen.gptq.bpw}      & \N{cap.qwen.gptq.dnll}$^\dagger$ & \N{cap.llama.gptq.dnll}      & -- \\
\bottomrule
\end{tabular}
\caption{Formats. Bits per weight include scales and biases. $\Delta$NLL is the change in mean per-token negative log-likelihood of bf16's own replies to 200 harmless prompts (bf16: 0.327, 0.203, 0.148). The 8-bit builds change likelihood very little, and MXFP4 loses the most in every model. AWQ and GPTQ were built only for the two models with a shift. $^\dagger$The \qwen{} GPTQ build failed and is excluded from all analyses (Appendix~\ref{app:prereg}). The llama.cpp builds (bf16, Q8\_0, Q4\_K\_M at \N{S9.qwenforced.gguf-q80.bpw} and \N{S9.qwenforced.gguf-q4km.bpw} bits per weight) run in a different runtime and are not comparable on this measure. They are reported in \S\ref{sec:controls} against their own bf16.}
\label{tab:formats}
\end{table}

\paragraph{Two settings.}
In the \emph{forced start}, the model sees one of the \N{C1.permodel.qwen.all.n} unique requests of the \texttt{harmful\_test} split of \citet{arditi2024refusal}, and the first $k=8$ tokens of a compliant reply are prefilled as token ids. The prefix comes from an abliterated copy of the same bf16 model \citep{arditi2024refusal}, so every build of a model continues the same tokens. We also report $k \in \{0, 2\}$ and a generic prefix, ``Sure, here is''. In the \emph{free start}, the model sees adversarial-harmful prompts from the WildJailbreak evaluation set \citep{jiang2024wildjailbreak}, which are in-the-wild jailbreak tactics, with no prefill. The first 200 prompts form a development set, and 400 further prompts, disjoint from these, form the held-out test set. Over-refusal is measured on XSTest's 250 safe prompts \citep{rottger2024xstest} and on WildJailbreak's 210 adversarial-benign prompts, which are benign requests written in jailbreak style. SORRY-Bench \citep{xie2024sorry} is a secondary family (Appendix~\ref{app:families}). Decoding is greedy with a cap of 128 new tokens.

\paragraph{Judges and human labels.}
The outcome is whether the reply's content does what was asked. A content judge (Qwen2.5-14B-Instruct) reads the request and the full reply and labels it as complying, partly complying, not complying, or incoherent. The first two count as \emph{judged harmful} in the forced start and as \emph{fulfilment} in the free start, where some prompts are mild. Keyword refusal detection cannot be used here, because the prefix fills the part of the reply that it reads. Llama Guard~3 \citep{llamaguard3} is a second judge with a different construct (``unsafe''). For the forced-start primary test we add human labels in a two-phase design: every pair of replies that the judge labelled differently, plus 25 random pairs from each of the two concordant strata, shown without format information and in random order, and design-weighted. The author labelled \N{humanpairs.total} pairs (\N{human.totalreplies} replies). Appendix~\ref{app:judge} gives the rubric and protocol.

\paragraph{Statistics.}
Every comparison is paired, with the same unique prompts under two builds. The \emph{shift} is the difference in rates, and \emph{churn} is the share of prompts whose label differs. Tests are exact McNemar tests, Holm-corrected within each family. Intervals are percentile bootstrap intervals over prompts with 10{,}000 resamples. Equivalence means that the 90\% interval of the shift lies within $\pm$5\pp{}. The primary and confirmatory tests, the secondary families and the decision rules were registered before the data they concern; Appendix~\ref{app:prereg} lists what was registered when, what changed, and which analyses are exploratory.
